\documentclass[11pt]{article}
\usepackage[margin=0.75in]{geometry}
\usepackage{amsmath,amssymb}
\usepackage{enumitem}
\usepackage{array}
\usepackage{booktabs}
\usepackage{colortbl}
\usepackage{xcolor}
\usepackage{tcolorbox}
\usepackage{tikz}
\usepackage{fancyhdr}
\usepackage{helvet}
\renewcommand{\familydefault}{\sfdefault}

% Colors
\definecolor{lightgray}{HTML}{F0F0F0}
\definecolor{medgray}{HTML}{CCCCCC}
\definecolor{darkrule}{HTML}{333333}
\definecolor{rubrichead}{HTML}{2C3E50}
\definecolor{rubriclight}{HTML}{EBF5FB}
\definecolor{rubricmid}{HTML}{D6EAF8}
\definecolor{frameblue}{HTML}{E8F4FD}
\definecolor{frameborder}{HTML}{3498DB}

% tcolorbox styles
\tcbset{
  wordbank/.style={
    colback=lightgray,
    colframe=medgray,
    boxrule=0.5pt,
    fonttitle=\bfseries\small,
    left=6pt, right=6pt, top=4pt, bottom=4pt,
    title=Word Bank
  },
  answerbox/.style={
    colback=lightgray!50,
    colframe=medgray,
    boxrule=0.5pt,
    left=6pt, right=6pt, top=6pt, bottom=6pt,
  },
  sentenceframe/.style={
    colback=frameblue,
    colframe=frameborder,
    boxrule=1pt,
    left=8pt, right=8pt, top=6pt, bottom=6pt,
    fonttitle=\bfseries,
    title=#1
  }
}

\newcommand{\blank}[1][2cm]{\underline{\hspace{#1}}}
\newcommand{\writeline}{\vspace{2pt}\noindent\rule{\linewidth}{0.3pt}\vspace{6pt}}

\pagestyle{empty}

\begin{document}

% ============================================================
%  PAGE 1: THE QUIZ
% ============================================================

\begin{center}
{\LARGE\bfseries Lesson 3-3 Quiz: Polynomial Identities}
\end{center}

\vspace{-6pt}
\noindent Name: \blank[7cm] \hfill Period: \blank[1.5cm]

\vspace{2pt}
\noindent\rule{\linewidth}{1.5pt}

\vspace{8pt}

% ---------- QUESTION 1 ----------
\noindent{\large\bfseries Question 1:} {\itshape Polynomial Identities (Factoring)} \hfill \textbf{/6 pts}

\vspace{6pt}
\noindent Use a polynomial identity to factor $\;8 + 125x^3$.

\vspace{8pt}
\noindent\textbf{(a)} Which polynomial identity applies here? Circle one. \hfill \textit{(1 pt)}

\vspace{4pt}
\noindent
\begin{tabular}{|p{0.46\textwidth}|p{0.46\textwidth}|}
\hline
\textbf{Difference of Cubes} & \textbf{Sum of Cubes} \\[2pt]
$a^3 - b^3 = (a - b)(a^2 + ab + b^2)$ & $a^3 + b^3 = (a + b)(a^2 - ab + b^2)$ \\[4pt]
\hline
\end{tabular}

\vspace{8pt}
\noindent\textbf{(b)} Show your work. Factor $\;8 + 125x^3\;$ completely. \hfill \textit{(3 pts)}

\vspace{4pt}
\begin{tcolorbox}[wordbank]
$2$ \qquad $5x$ \qquad $4$ \qquad $25x^2$ \qquad $10x$ \qquad sum of cubes \qquad difference of cubes
\end{tcolorbox}

\vspace{6pt}
\noindent Identify $a$ and $b$: \qquad $a = $ \blank[2cm] \qquad $b = $ \blank[2cm]

\vspace{10pt}
\noindent Substitute into the identity:

\vspace{24pt}

\noindent Final factored form: \quad $8 + 125x^3 \;=\; \bigl($\blank[2cm]$\bigr)\bigl($\blank[3.5cm]$\bigr)$

\vspace{10pt}
\noindent\textbf{(c)} In 1--2 sentences, explain how you knew which identity to use. \hfill \textit{(2 pts)}

\vspace{2pt}
\writeline
\writeline

\vspace{2pt}
\noindent\rule{\linewidth}{1.5pt}
\vspace{6pt}

% ---------- QUESTION 2 ----------
\noindent{\large\bfseries Question 2:} {\itshape Binomial Theorem Expansion} \hfill \textbf{/4 pts}

\vspace{6pt}
\noindent Use the Binomial Theorem to expand $\;(x - 3)^5$.

\vspace{4pt}
\begin{tcolorbox}[wordbank]
Pascal's Triangle coefficients for $n = 5$: \quad 1 \qquad 5 \qquad 10 \qquad 10 \qquad 5 \qquad 1
\end{tcolorbox}

\vspace{6pt}
\noindent\textbf{(a)} Fill in the missing coefficients from Pascal's Triangle: \hfill \textit{(1 pt)}

\vspace{4pt}
\noindent\hspace{6pt}$\displaystyle(x-3)^5 \;=\; $\blank[0.6cm]$x^5(-3)^0 \;+\;$\blank[0.6cm]$x^4(-3)^1 \;+\;$\blank[0.6cm]$x^3(-3)^2$

\vspace{6pt}
\noindent\hspace{62pt}$+\;$\blank[0.6cm]$x^2(-3)^3 \;+\;$\blank[0.6cm]$x^1(-3)^4 \;+\;$\blank[0.6cm]$(-3)^5$

\vspace{10pt}
\noindent\textbf{(b)} Evaluate each power of $(-3)$ and simplify each term: \hfill \textit{(2 pts)}

\vspace{24pt}

\noindent\textbf{(c)} Write the expanded polynomial in standard form: \hfill \textit{(1 pt)}

\vspace{4pt}
\begin{tcolorbox}[answerbox]
\textbf{Final Answer:} \quad $(x - 3)^5 = $ \blank[10cm]
\end{tcolorbox}

\vspace{6pt}
\hfill \textbf{Total: \qquad /10}

\newpage

% ============================================================
%  PAGE 2: SCORING RUBRIC
% ============================================================

\begin{center}
{\LARGE\bfseries Lesson 3-3 Quiz --- Scoring Rubric}\\[4pt]
{\small\itshape Teacher Copy \quad$\cdot$\quad 10 points total}
\end{center}

\vspace{4pt}
\noindent\rule{\linewidth}{1.5pt}
\vspace{8pt}

% ----- Q1 Rubric -----
{\large\bfseries Question 1: Polynomial Identity Factoring \hfill 6 points}

\vspace{6pt}

\renewcommand{\arraystretch}{1.4}
\noindent
\begin{tabular}{>{\bfseries}p{1.2cm} p{5.5cm} p{3cm} p{4.5cm}}
\toprule
\rowcolor{rubrichead}\color{white}Part & \color{white}Criteria & \color{white}Points & \color{white}Answer Key \\
\midrule
\rowcolor{rubriclight}
(a) & Correctly identifies \textbf{Sum of Cubes} & 1 pt & Sum of Cubes \\
\rowcolor{white}
(b) & Correctly identifies $a = 2$ and $b = 5x$ & 1 pt & $a=2,\; b=5x$ \\
\rowcolor{rubriclight}
(b) & Shows substitution into identity & 1 pt & $(2+5x)(4-10x+25x^2)$ \\
\rowcolor{white}
(b) & Correct final factored form & 1 pt & $(2+5x)(4-10x+25x^2)$ \\
\rowcolor{rubriclight}
(c) & States the expression is a \textit{sum} (uses $+$) & 1 pt & \textit{See sample below} \\
\rowcolor{white}
(c) & Identifies both terms as perfect cubes ($2^3=8$, $(5x)^3=125x^3$) & 1 pt & \textit{See sample below} \\
\bottomrule
\end{tabular}

\vspace{8pt}
\begin{tcolorbox}[colback=white, colframe=rubrichead, title={\color{white}Sample response for (c)}, fonttitle=\bfseries\small, coltitle=white]
\textit{``I knew to use the sum of cubes identity because the expression uses addition ($+$), and both 8 and $125x^3$ are perfect cubes: $2^3 = 8$ and $(5x)^3 = 125x^3$.''}
\end{tcolorbox}

\vspace{12pt}

% ----- Q2 Rubric -----
{\large\bfseries Question 2: Binomial Theorem Expansion \hfill 4 points}

\vspace{6pt}

\noindent
\begin{tabular}{>{\bfseries}p{1.2cm} p{6cm} p{1.5cm} p{5cm}}
\toprule
\rowcolor{rubrichead}\color{white}Part & \color{white}Criteria & \color{white}Points & \color{white}Answer Key \\
\midrule
\rowcolor{rubriclight}
(a) & Fills in all 6 Pascal's Triangle coefficients correctly: 1, 5, 10, 10, 5, 1 & 1 pt & Coefficients provided in word bank \\
\rowcolor{white}
(b) & Correctly evaluates powers of $(-3)$:\newline $1,\; -3,\; 9,\; -27,\; 81,\; -243$ & 1 pt & Common error: all positive (ignoring negative sign) \\
\rowcolor{rubriclight}
(b) & Correctly multiplies coefficients $\times$ powers of $(-3)$ to simplify each term & 1 pt & $1,\; -15,\; 90,\; -270,\; 405,\; -243$ \\
\rowcolor{white}
(c) & Writes complete polynomial in standard form & 1 pt & See key below \\
\bottomrule
\end{tabular}

\vspace{8pt}
\begin{tcolorbox}[colback=white, colframe=rubrichead, title={\color{white}Answer Key}, fonttitle=\bfseries\small, coltitle=white]
\vspace{-4pt}
\begin{align*}
(x-3)^5 &= \binom{5}{0}x^5(-3)^0 + \binom{5}{1}x^4(-3)^1 + \binom{5}{2}x^3(-3)^2 + \binom{5}{3}x^2(-3)^3 + \binom{5}{4}x^1(-3)^4 + \binom{5}{5}x^0(-3)^5 \\[4pt]
&= x^5 - 15x^4 + 90x^3 - 270x^2 + 405x - 243
\end{align*}
\end{tcolorbox}

\vspace{10pt}
\noindent\rule{\linewidth}{0.5pt}
\vspace{6pt}

{\bfseries Topic 3 Assessment Alignment:}\\[4pt]
\begin{tabular}{ll}
Q1 (Sum of Cubes) & $\longrightarrow$ Topic Assessment Q8: Factor $64 + 27a^3$ \\
Q2 (Binomial Expansion) & $\longrightarrow$ Topic Assessment Q3: Expand $(x-2)^6$ and Q13: Expand $(2a+2b)^5$ \\
\end{tabular}

\newpage

% ============================================================
%  PAGE 3: SENTENCE FRAMES (ELL SUPPORT)
% ============================================================

\begin{center}
{\LARGE\bfseries Lesson 3-3 Quiz --- Sentence Frames}\\[4pt]
{\small\itshape ELL Support Version \quad$\cdot$\quad Use with Question 1, Part (c)}
\end{center}

\vspace{4pt}
\noindent\rule{\linewidth}{1.5pt}
\vspace{10pt}

\noindent{\large\bfseries Question 1(c):} In 1--2 sentences, explain how you knew which identity to use.

\vspace{12pt}

% ----- FRAME OPTION A -----
\begin{tcolorbox}[sentenceframe={Option A --- Fill in the Blanks}]
\vspace{4pt}

I used the \blank[3.5cm] identity because the two terms are connected by \blank[2.5cm]\,.

\vspace{14pt}
I know that \blank[1.5cm] is a perfect cube because \blank[1.5cm]$\,^3 = $ \blank[1.5cm]\,, 

\vspace{14pt}
and \blank[2cm] is a perfect cube because \blank[1.5cm]$\,^3 = $ \blank[2cm]\,.

\vspace{6pt}
\end{tcolorbox}

\vspace{6pt}

\begin{tcolorbox}[wordbank]
$8$ \qquad $125x^3$ \qquad $2$ \qquad $5x$ \qquad addition ($+$) \qquad subtraction ($-$)\\[4pt]
sum of cubes \qquad difference of cubes
\end{tcolorbox}

\vspace{16pt}

% ----- FRAME OPTION B -----
\begin{tcolorbox}[sentenceframe={Option B --- Sentence Starters}]
\vspace{4pt}

I chose the \blank[5cm] identity because the expression $8 + 125x^3$ shows \blank[5cm]\,.

\vspace{14pt}
Both terms are perfect cubes because \rule{0pt}{12pt}\hrulefill

\vspace{14pt}
\hrulefill

\vspace{6pt}
\end{tcolorbox}

\vspace{20pt}

% ----- KEY VOCABULARY -----
\begin{tcolorbox}[colback=white, colframe=darkrule, boxrule=1pt, title={\color{white}\bfseries Key Vocabulary}, fonttitle=\large, coltitle=white, colbacktitle=darkrule, left=8pt, right=8pt, top=6pt, bottom=6pt]
\renewcommand{\arraystretch}{1.5}
\begin{tabular}{>{\bfseries}p{3cm} p{11.5cm}}
identity & A math equation that is always true for all values of the variable. \\
perfect cube & A number or expression that can be written as something raised to the 3rd power.\\
 & \textit{Examples:} $8 = 2^3$ \qquad $125x^3 = (5x)^3$ \qquad $27 = 3^3$ \\
sum & The result of \textbf{adding}. The symbol $+$ means sum. \\
difference & The result of \textbf{subtracting}. The symbol $-$ means difference. \\
factor & To rewrite an expression as a product (multiplication) of simpler parts. \\
\end{tabular}
\end{tcolorbox}

\end{document}
